% !TeX program = pdflatex
% =============================================================================
% Arbeitsblatt 4: Normalkraft
% UZH Praktikum I -- Sara Romer Urban, Kantonsschule im Lee, HS2026
% Klasse: 2e (09.09.2026)
%
% Folgt dem in normalkraft-lektionsplan.tex festgelegten Ablauf: Repetition
% (Feder-Gleichgewicht) -> Tisch/kognitiver Konflikt (Bilder/Normalkraft_Tisch.png)
% -> elastische Unterlage als sichtbarer Beleg (Bilder/Normalkraft_elastische_
% oberflaeche.png) -> Definition Normalkraft, mit Tisch-Bild erneut zur
% Generalisierung (auch scheinbar starre Oberflaechen federn mikroskopisch) und
% mikroskopischer Ursache (Bilder/atoms-repelling-normal-force.jpg) -> Klotz-
% Sortieraufgabe (LZ2, Bilder/Normalkraft_klotz(.png/_loesung.png)) -> schiefe
% Ebene mit der Bergziege bei drei Winkeln (LZ3/LZ4,
% Bilder/schiefe_ebene_1.png, Loesungen schiefe_ebene_1_loesung_
% kraeftezerlegung.png/_normalkraft.png) -> Anwendungsaufgabe Endgeschwindigkeit
% (LZ5, Werte aus lernziele.md: m=80kg, s=10m, F_HA~390N, a~4.9m/s^2, v~9.9m/s).
%
% Bewusst nicht erwaehnt (siehe lernziele.md "Bewusst nicht Teil"): Aktio-Reaktio/
% drittes Newtonsches Gesetz, das Koerper-auf-Unterlage-Gegenstueck. Die Bilder
% Normalkraft_Tisch.png/_elastische_oberflaeche.png zeigen zwar auch eine
% -F_N-Beschriftung auf der Unterlagenseite, der Fliesstext hier benennt aber
% ausschliesslich F_N auf den Koerper -- die Wechselwirkungs-Lektion (16.09.2026)
% bleibt fuer die Interaktionspaar-Perspektive zustaendig.
%
% Praktikum I: Sie-Form, keine farbigen Boxen ausser Formel-/Hinweis-/Antwortbox,
% Textmenge knapp (Sara Romers Feedback ab Kondensatoren). Boxen-Definitionen aus
% arbeitsblatt3-kraefte-addieren-zerlegen.tex uebernommen.
%
% Quellen: lernziele.md, normalkraft-lektionsplan.tex (dieser Ordner),
% Bilder/Normalkraft_Tisch.png, Bilder/Normalkraft_elastische_oberflaeche.png,
% Bilder/atoms-repelling-normal-force.jpg, Bilder/Normalkraft_klotz(.png/
% _loesung.png), Bilder/schiefe_ebene_1.png, Bilder/schiefe_ebene_1_loesung_
% kraeftezerlegung.png, Bilder/schiefe_ebene_1_loesung_normalkraft.png.
% =============================================================================
\documentclass[addpoints,11pt,a4paper]{exam}

\usepackage[T1]{fontenc}
\usepackage[utf8]{inputenc}
\usepackage[ngerman]{babel}
\usepackage{lmodern}
\usepackage[a4paper,margin=2.2cm,headheight=14pt]{geometry}
\usepackage{amsmath}
\usepackage{siunitx}
% Hausstil: zusammengesetzte Einheiten immer als echter Bruch (\frac), nicht
% als Inline-Schraegstrich oder \tfrac -- gilt global fuer jedes \si/\SI.
\sisetup{per-mode=fraction,fraction-command=\frac}
\usepackage{array,booktabs,tabularx,multirow}
\usepackage{enumitem}
\usepackage{xcolor}
\usepackage{colortbl}
\usepackage{tikz}
\usepackage{physics}
\usepackage[most]{tcolorbox}
\usepackage{lastpage}
\usepackage{graphicx}
\usepackage{hyperref}

\usetikzlibrary{arrows.meta,calc,angles,quotes,decorations.pathmorphing}

\usepackage{setspace}
\usepackage{needspace}
\setlength{\parindent}{0pt}
\setlength{\parskip}{0.6em}
\setstretch{1.08}
\setlist[itemize]{itemsep=0.4em}

% -----------------------------
% Farben (Corporate-Farbschema Physik KFR)
% -----------------------------
\definecolor{titleblue}{HTML}{183A6B}
\definecolor{accentblue}{HTML}{2E6F9E}
\definecolor{lightblue}{HTML}{EAF4FB}
\definecolor{lightgray}{HTML}{F4F6F8}
\definecolor{darkgray}{HTML}{4A4A4A}
\definecolor{accentorange}{HTML}{D9720C}
\definecolor{accentpurple}{HTML}{7B2D8E}
% Hinweisbox: Orange (accentorange), fuer wichtige Klarstellungen.
\definecolor{lightorange}{HTML}{FCEEE0}
% Formelbox: Violett (accentpurple), fuer die zentralen Merksaetze.
\definecolor{lightpurple}{HTML}{F3EAF7}

% Links (hyperref): farblich ans Corporate-Farbschema angeglichen.
\hypersetup{
  colorlinks=true,
  urlcolor=accentblue,
  linkcolor=accentblue,
  pdftitle={Arbeitsblatt 4: Normalkraft},
  pdfauthor={Loris Keller}
}

% -----------------------------
% Spaltentypen
% -----------------------------
\newcolumntype{Y}{>{\centering\arraybackslash}X}
\newcolumntype{P}[1]{>{\raggedright\arraybackslash}p{#1}}

% -----------------------------
% Boxen
% -----------------------------
\tcbset{before skip=10pt, after skip=10pt}
\newtcolorbox{titlebox}{
  enhanced,
  colback=titleblue,
  colframe=titleblue,
  boxrule=0pt,
  arc=3mm,
  left=3mm,right=3mm,top=3mm,bottom=3mm
}

\newlength{\aufgabenindent}
\settowidth{\aufgabenindent}{10.\hskip\labelsep}

% Praktikum I: Lernziele/Einleitung/Theorie/Aufgaben werden NICHT mehr
% farbig gerahmt (nur Formel-, Hinweis- und Antwortbox bleiben Boxen) --
% siehe institutionen/UZH-Praktikum-I/CLAUDE.md, Abschnitt "Boxen-Layout".
\newenvironment{infobox}{%
  \par\addvspace{10pt}\noindent
}{%
  \par\addvspace{10pt}
}

\newenvironment{theoriebox}[1][Theorie]{%
  \par\addvspace{10pt}\noindent\textbf{#1}\par\smallskip\noindent
}{%
  \par\addvspace{10pt}
}

% taskbox/groupbox stehen in questions VOR dem ersten \question (Bilder,
% Fliesstext) -- exam.cls' questions-Liste braucht dafuer eine echte,
% eigenstaendige Box (sonst "Something's wrong--perhaps a missing \item",
% weil z.B. ein \begin{center} vor dem ersten \item der Aussenliste steht).
% Deshalb bleiben es tcolorboxen, nur unsichtbar gemacht (keine Farbe/Rahmen,
% kein Padding) statt sie durch reine Absaetze zu ersetzen.
\newtcolorbox{taskbox}[1][]{
  enhanced, breakable,
  colback=white, colframe=white, boxrule=0pt,
  left=0mm,right=0mm,top=0mm,bottom=0mm,
  before skip=10pt, after skip=10pt,
  before upper={%
    \def\tmpboxtitle{#1}%
    \ifx\tmpboxtitle\empty\else\noindent\textbf{#1}\par\smallskip\fi
  }
}

\newtcolorbox{groupbox}[1][Gruppenarbeit]{
  enhanced, breakable,
  colback=white, colframe=white, boxrule=0pt,
  left=0mm,right=0mm,top=0mm,bottom=0mm,
  before skip=10pt, after skip=10pt,
  before upper={\noindent\textbf{#1}\par\smallskip}
}

% beispielbox: fuer die Einleitung -- wie die anderen Inhaltsboxen ungerahmt.
\newenvironment{beispielbox}[1][Beispiel]{%
  \par\addvspace{10pt}\noindent\textbf{#1}\par\smallskip\noindent
}{%
  \par\addvspace{10pt}
}

% hinweisbox: Orange, fuer wichtige Klarstellungen/Verwechslungswarnungen.
\newtcolorbox{hinweisbox}[1][Hinweis]{
  enhanced,
  breakable,
  colback=lightorange,
  colframe=accentorange,
  title={#1},
  fonttitle=\bfseries,
  boxrule=0.7pt,
  arc=2mm,
  left=5mm,right=5mm,top=2mm,bottom=2mm
}

% formelbox: Violett, um die zentralen Merksaetze dieser Einheit hervorzuheben.
\newtcolorbox{formelbox}[1][Merksatz]{
  enhanced,
  breakable,
  colback=lightpurple,
  colframe=accentpurple,
  title={#1},
  fonttitle=\bfseries,
  boxrule=0.9pt,
  arc=2mm,
  left=5mm,right=5mm,top=2mm,bottom=2mm
}

\newtcolorbox{answerbox}{
  breakable,
  colback=white,
  colframe=gray!55,
  boxrule=0.5pt,
  arc=1.5mm,
  left=2mm,right=2mm,top=2mm,bottom=2mm
}

% Marke: kleines farbiges Inline-Label fuer Teilschritte/Ergebnis-Callouts.
\newtcbox{\Marke}[1][accentpurple]{
  nobeforeafter,
  colback=#1,
  colframe=#1,
  coltext=white,
  boxrule=0pt,
  arc=2pt,
  left=5pt,right=5pt,top=2pt,bottom=2pt,
  fontupper=\bfseries\sffamily\small
}

% Bild-Helfer: zentriertes Bild mit spuerbarem Abstand zum umgebenden Text
% (statt nur des normalen Absatzabstands). #1 = optionale includegraphics-
% Optionen (z. B. trim/clip, siehe schiefe_ebene_1*-Bilder unten), #2 =
% Breite, #3 = Dateipfad.
\newcommand{\Bild}[3][]{%
  \par\vspace{0.6cm}
  \begin{center}
    \includegraphics[#1,width=#2]{#3}
  \end{center}
  \vspace{0.6cm}
}

% --- Loesungen ein-/ausblenden -----------------------------------------------
%\noprintanswers
\printanswers

\pointformat{}
\renewcommand{\solutiontitle}{\noindent\color{titleblue}\textbf{L\"osung:}\enspace}

\pagestyle{headandfoot}
\cfoot{\textcolor{darkgray}{Seite \thepage\ von \pageref{LastPage}}}

\newcommand{\Kopfzeile}[4]{%
  \begin{titlebox}
    {\color{white}\sffamily\bfseries\Large #4}
  \end{titlebox}
  \vspace{0.5em}
  \noindent\textbf{Klasse:}~#1 \hfill \textbf{Datum:}~#2 \hfill \textbf{Thema:}~#3
    \hfill \textbf{Lehrperson:}~Loris Keller
  \vspace{0.8em}
  \lhead{\textcolor{darkgray}{#3}}
  \rhead{\textcolor{darkgray}{#4}}
}

\newlength{\antwortraumlen}
\newcommand{\Antwortraum}[1]{%
  \pgfmathsetlength{\antwortraumlen}{#1*2.5cm}%
  \begin{answerbox}
    \rule{0pt}{\antwortraumlen}%
  \end{answerbox}%
}

% Werte fuer Kopfzeile zentral setzen
\newcommand{\Klasse}{2e}
\newcommand{\Datum}{09.09.2026}
\newcommand{\Thema}{Normalkraft}
\newcommand{\ArbeitsblattTitel}{Arbeitsblatt 4: Normalkraft}

\begin{document}

\Kopfzeile{\Klasse}{\Datum}{\Thema}{\ArbeitsblattTitel}

% =============================================================================
% Repetition: Kraeftegleichgewicht Feder-Masse
% =============================================================================
\textbf{Repetition}

\begin{center}
\begin{tikzpicture}[>=Latex,scale=1.3]
  \draw[very thick] (-1,3) -- (1,3);
  \foreach \x in {-0.8,-0.4,...,0.8} { \draw (\x,3) -- (\x-0.25,3.3); }
  \draw[decorate,decoration={coil,aspect=0.3,segment length=3.2pt,amplitude=3pt}] (0,3) -- (0,1.8);
  \draw[fill=lightgray!60] (-0.4,1.3) rectangle (0.4,1.8);
  \node at (0,1.55) {$m$};
  \fill (0,1.55) circle (0.5pt);
  \draw[->,ultra thick,titleblue] (0,1.55) -- ++(0,1.1) node[midway,right]{$\vec F$};
  \draw[->,ultra thick,accentorange] (0,1.55) -- ++(0,-1.1) node[midway,right]{$\vec F_G$};
\end{tikzpicture}
\end{center}

Eine Federwaage hält eine Masse $m$ in Ruhe: Federkraft und
Gewichtskraft befinden sich im Kräftegleichgewicht,
$\vec F+\vec F_G=\vec 0$.

% =============================================================================
% Theorie: Definition Normalkraft (Tisch + elastische Unterlage als Belege)
% =============================================================================
\begin{theoriebox}[Normalkraft]
Legt man dieselbe Masse stattdessen auf einen Tisch, bewegt sie sich
trotzdem nicht durch den Tisch hindurch, obwohl nur die Gewichtskraft
$\vec F_G$ sichtbar nach unten zieht:

\Bild{7cm}{Bilder/Normalkraft_Tisch.png}

Der Tisch muss also eine zweite, unsichtbare Kraft ausüben. Auf einer
weichen, nachgiebigen Unterlage sieht man diese Gegenkraft direkt: Die
Unterlage verformt sich und drückt den Körper dadurch sichtbar nach
oben zurück.

\Bild{7cm}{Bilder/Normalkraft_elastische_oberfläche.png}

Jede Oberfläche, auf der ein Körper aufliegt, übt als Reaktion darauf
eine Gegenkraft aus: die \textbf{\emph{Normalkraft}} $F_N$ (Einheit:
$\si{\newton}$). Sie ist eine \textbf{\emph{Reaktionskraft}} und zeigt
immer in die Richtung \textbf{\emph{senkrecht zur Oberfläche}}. Sie
verhindert, dass der Körper in die Oberfläche eindringt. Beim starren
Tisch von vorhin gilt genau dasselbe Prinzip, nur verformt er sich dabei
mikroskopisch wenig, für das Auge unsichtbar.

\textbf{\emph{Mikroskopische Ursache:}} Oberflächen bestehen aus Atomen/
Molekülen, die sich bei Annäherung elektrisch abstossen, wie an einem
Federmodell sichtbar wird, wenn eine Hand auf ein Atomgitter drückt:

\Bild{6cm}{Bilder/atoms-repelling-normal-force.jpg}
\end{theoriebox}

% =============================================================================
% Uebung: Klotz-Sortieraufgabe (LZ2)
% =============================================================================
\needspace{6cm}
\begin{questions}
\begin{taskbox}[Aufgabe: Normalkraft oder Gewichtskraft: was ist grösser?]
\question[4] Vier Situationen mit demselben Klotz (Gewichtskraft
$\vec F_g$ immer gleich gross):

\Bild{13cm}{Bilder/Normalkraft_klotz.png}

Ordnen Sie für jede Situation $F_N$ und $F_g$ zueinander ein: kleiner,
gleich gross oder grösser?

\Antwortraum{1.5}

\begin{solution}
\Bild{13cm}{Bilder/Normalkraft_klotz_loesung.png}
\end{solution}
\end{taskbox}
\end{questions}

% =============================================================================
% Uebergang und Theorie: Schiefe Ebene
% =============================================================================
Wie in der letzten Lektion zerlegen wir Kräfte grafisch mit dem
Kräfteparallelogramm, jetzt für eine geneigte Oberfläche.

\begin{theoriebox}[Schiefe Ebene: Gewichtskraft in zwei Komponenten zerlegen]
Auf einer geneigten Ebene zeigt die Gewichtskraft $\vec F_G$ weiterhin
senkrecht nach unten, die Normalkraft $\vec F_N$ aber senkrecht zur
\emph{geneigten} Ebene. Deshalb zerlegen wir $\vec F_G$ grafisch in zwei
Komponenten:

\begin{itemize}[leftmargin=1.2cm,itemsep=0.2em]
  \item die \textbf{\emph{parallele Komponente}} $F_{G\parallel}$
  (Einheit: $\si{\newton}$), parallel zur Ebene,
  \item die \textbf{\emph{senkrechte Komponente}} $F_{G\perp}$
  (ebenfalls in $\si{\newton}$), senkrecht zur Ebene, wobei ihr die
  Normalkraft $F_N$ exakt das Gleichgewicht hält.
\end{itemize}

Konstruktion: Hilfslinien parallel und senkrecht zur Ebene durch die
Spitze von $\vec F_G$ ziehen. Weil diese Hilfslinien senkrecht
aufeinander stehen, wird aus dem Kräfteparallelogramm ein Rechteck mit
$\vec F_G$ als Diagonale. Beide Rechtecksseiten mit dem Kraftmassstab
ausmessen.
\end{theoriebox}

% =============================================================================
% Uebung: Konstruktion fuer drei Neigungswinkel (LZ3, LZ4)
% =============================================================================
\begin{questions}
\setcounter{question}{1}
\begin{taskbox}[Aufgabe: Die Bergziege auf der vereisten Bergwand]
\Bild{7cm}{Bilder/bergziege_damm.jpg}

\question[9] Dieselbe Bergziege steht bei drei unterschiedlichen
Neigungswinkeln $\alpha_1$, $\alpha_2$, $\alpha_3$:

% trim/clip: schneidet einen fehlerhaften grauen Streifen oben in der
% Bildquelle weg (Exportartefakt, alle drei schiefe_ebene_1*-Bilder
% betroffen, Pixel y=0-169 von 2845 bei allen dreien identisch).
\Bild[trim=0 0 0 50bp,clip]{\textwidth}{Bilder/schiefe_ebene_1.png}

Konstruieren Sie für $\alpha_1$ gemeinsam mit der Lehrperson, für
$\alpha_2$ und $\alpha_3$ selbständig direkt in der Skizze: die
parallele Komponente $\vec F_{G\parallel}$, die senkrechte Komponente
$\vec F_{G\perp}$ und die Normalkraft $\vec F_N$. Vergleichen Sie
anschliessend, wie sich $F_{G\parallel}$ und $F_N$ mit zunehmender
Steigung verändern.

\begin{solution}
\Bild[trim=0 0 0 50bp,clip]{\textwidth}{Bilder/schiefe_ebene_1_loesung_kräftezerlegung.png}
\Bild[trim=0 0 0 50bp,clip]{\textwidth}{Bilder/schiefe_ebene_1_loesung_normalkraft.png}

Je steiler die Ebene, desto grösser $F_{G\parallel}$ und desto kleiner
$F_N$.
\end{solution}
\end{taskbox}
\end{questions}

% =============================================================================
% Anwendung: Endgeschwindigkeit (LZ5)
% =============================================================================
\begin{questions}
\setcounter{question}{2}
\begin{taskbox}[Aufgabe: Die Bergziege rutscht ab]
\question[6] Bei $\alpha_2$ verliert die Bergziege (Masse
$m=\SI{80}{\kilogram}$) auf dem reibungsfreien Eis den Halt und rutscht
aus der Ruhe eine Hanglänge von $s=\SI{10}{\meter}$ hinunter.

\begin{parts}
\part[3] Welche der beiden Kräfte aus Ihrer Konstruktion ist für die
Beschleunigung entlang des Hangs verantwortlich? Nutzen Sie diese Kraft
und das 2. Newtonsche Gesetz, um die Beschleunigung $a$ zu berechnen.

\Antwortraum{2}

\begin{solution}
Verantwortlich ist die parallele Komponente $F_{G\parallel}$ der
Gewichtskraft, nicht die volle Gewichtskraft $F_g$: Die Normalkraft
$F_N$ hält bereits der senkrechten Komponente $F_{G\perp}$ das
Gleichgewicht. Erwartungsgemäss $F_{G\parallel}\approx\SI{390}{\newton}$.
Nach $a$ auflösen:
\[
F_{G\parallel}=m\cdot a \quad\Longrightarrow\quad a=\frac{F_{G\parallel}}{m}
\]
Werte einsetzen:
\[
a=\frac{\SI{390}{\newton}}{\SI{80}{\kilogram}}\approx\SI{4.9}{\meter\per\second\squared}
\]
\end{solution}

\part[3] Bestimmen Sie mit $v^2=2\cdot a\cdot s$ die Endgeschwindigkeit
$v$ der Bergziege am Ende des Hangs.

\Antwortraum{2}

\begin{solution}
Nach $v$ auflösen:
\[
v^2=2\cdot a\cdot s \quad\Longrightarrow\quad v=\sqrt{2\cdot a\cdot s}
\]
Werte einsetzen:
\[
v=\sqrt{2\cdot\SI{4.9}{\meter\per\second\squared}\cdot\SI{10}{\meter}}
  \approx\SI{9.9}{\meter\per\second}
\]
\end{solution}
\end{parts}
\end{taskbox}
\end{questions}

% =============================================================================
% Abschluss
% =============================================================================
Weitere Übungsaufgaben zur Normalkraft finden Sie im separaten
Aufgabenblatt.

\end{document}
